\chapter{Evolução dos CLPs}
\label{cap:evolucao}

\begin{objetivos}
Ao final deste capítulo você deve ser capaz de:
\begin{itemize}
  \item descrever o problema industrial que deu origem ao CLP;
  \item situar no tempo as gerações de controladores e o que mudou em cada uma;
  \item diferenciar CLP, PAC, controlador de segurança e controlador virtual;
  \item reconhecer o que permaneceu igual desde 1969 e por quê.
\end{itemize}
\end{objetivos}

\section{Antes do CLP: a lógica fiada}

O controle de máquinas nasceu eletromecânico. Relés, contatos, bobinas e temporizadores de motor
síncrono executavam a sequência, e a lógica estava fisicamente desenhada na fiação. Funcionava
bem --- mas cobrava um preço alto a cada mudança:

\begin{itemize}
  \item mau contato e desgaste mecânico dos contatos, com falhas intermitentes;
  \item instalação de um número enorme de relés, com fiação ponto a ponto entre centenas de
        terminais;
  \item dificuldade de alterar a sequência, que exigia redesenhar e refazer o painel;
  \item manutenção periódica obrigatória e paradas recorrentes.
\end{itemize}

Mudar a lógica de uma linha de montagem significava parar a produção e trocar os quadros
elétricos. Era gasto de tempo e de dinheiro em uma escala que a indústria automobilística não
podia mais aceitar.

\index{relé!painel de comando}\index{lógica fiada}

\section{1968--1969: a especificação da GM e o primeiro CLP}

Em \textbf{1968}, a divisão \emph{Hydramatic} da General Motors --- responsável pelas transmissões
automáticas --- publicou uma especificação para um substituto eletrônico dos painéis de relé. A
especificação nasceu de um documento do engenheiro Edward R. Clark e refletia a insatisfação de
muitos usuários de relés, dentro e fora da indústria automobilística. Na época, a liderança do
grupo ficou associada ao engenheiro \textbf{Richard ``Dick'' Morley}, hoje considerado o pai do CLP.

A proposta vencedora foi da \textbf{Bedford Associates}, de Bedford, Massachusetts. O equipamento
ficou pronto em \textbf{1969} e recebeu a designação \textbf{084}, porque era o octogésimo quarto
projeto da empresa. A Bedford criou em seguida a \textbf{Modicon}, dedicada a fabricar e vender o
produto --- o nome vinha de \emph{modular digital controller}. Em paralelo, a Allen-Bradley
desenvolvia sua própria linha, e o engenheiro Odo J. Struger ficou conhecido por cunhar a sigla
\emph{PLC} para distinguir o controlador do ``PC'' (\emph{personal computer}).\footnote{Fatos
históricos conferidos em Laughton \& Warne (2002) e Dunn (2009), reunidos em
\url{https://en.wikipedia.org/wiki/Programmable_logic_controller} (consulta em 22/09/2026).}

A Modicon criou também o \textbf{Modbus}, protocolo que se tornaria padrão aberto
(Capítulo~\ref{cap:redes-scada}).

\section{Os dez requisitos da especificação original}

A especificação de 1968 é um excelente roteiro de projeto porque continua atual. O novo
equipamento deveria:

\begin{table}[H]
  \centering
  \caption{Requisitos da especificação original, reescritos e ainda válidos.}
  \label{tab:requisitos-originais}
  \begin{tabular}{p{0.9cm} p{6.0cm} p{7.4cm}}
    \toprule
    \# & Requisito original & Leitura atual \\
    \midrule
    1 & Facilidade de programar e de alterar o programa, inclusive na estação de trabalho
      & Programação em campo, com o processo rodando, e alteração sem refazer painel \\
    2 & Manutenção fácil, de preferência com unidades \emph{plug-in}
      & Módulos plug-in e troca a quente \\
    3 & Confiabilidade superior à do painel de relés
      & É a razão de existir do CLP: ambiente industrial, imunidade a ruído, diagnóstico \\
    4 & Mais compacto que o painel de relés
      & Densidade de pontos por módulo e por rack \\
    5 & Envio direto de dados à unidade central de processamento
      & O embrião do que hoje é a integração com SCADA e IIoT (Capítulos~\ref{cap:iiot}--\ref{cap:dados}) \\
    6 & Custo competitivo frente ao painel de relés
      & Justificativa econômica do projeto; hoje inclui custo de ciclo de vida \\
    7 & Aceitar entradas em 115\,V CA
      & Compatibilidade com o sinal de campo da instalação: 24\,V CC e 4--20\,mA hoje \\
    8 & Fornecer saídas em 115\,V CA acima de 2\,A, capazes de acionar válvulas e partidas de motor
      & Acionamento direto de cargas com a saída discreta \\
    9 & Permitir ampliação do sistema com o mínimo de alteração
      & Arquitetura modular e expansão por módulos e remotas \\
    10 & Ser programável por pessoal de manutenção, não só por especialistas
      & Linguagem gráfica próxima do esquema elétrico: a razão do sucesso do Ladder \\
    \bottomrule
  \end{tabular}
\end{table}

\section{Gerações: o que mudou em cada salto}

\begin{figure}[H]
  \centering
  \includegraphics[width=\linewidth]{02-linha-do-tempo.pdf}
  \caption{Linha do tempo da evolução dos controladores industriais.}
  \label{fig:linha-do-tempo}
\end{figure}

\begin{table}[H]
  \centering
  \caption{Gerações de controladores e o que mudou em cada uma.}
  \label{tab:geracoes}
  \begin{tabular}{p{2.1cm} p{4.1cm} p{4.1cm} p{4.3cm}}
    \toprule
    Período & Programação & Hardware e memória & Limitação da época \\
    \midrule
    1969--1975 & Terminal dedicado, lógica booleana, sem representação gráfica; depois Ladder
      & Memória de núcleo magnético, componentes discretos; programa em fita cassete
      & Sem padronização; pouca capacidade de documentação e impressão \\
    1975--1985 & Ladder consolidado como linguagem de campo
      & Microprocessador substitui a lógica discreta; memória semicondutora
      & Redes proprietárias e fechadas; endereçamento específico de cada fabricante \\
    1985--1995 & Primeiros softwares de programação em PC; sub-rotinas
      & CLP modular em rack, maior densidade de pontos
      & Programas presos ao fabricante; troca de marca exigia reescrever tudo \\
    1995--2005 & Padronização das linguagens pela IEC 61131-3
      & Barramentos de campo, I/O remoto, CPUs mais rápidas
      & Portabilidade ainda parcial: endereçamento e blocos de sistema continuam proprietários \\
    2005--2015 & Ambiente integrado de projeto, simulação e versionamento
      & Ethernet industrial, controle distribuído, redundância
      & Integração com a rede corporativa traz o problema de segurança cibernética \\
    2015--2026 & Padrões de modelagem de informação (OPC UA) e integração com software livre
      & Borda (\emph{edge}), virtualização e controladores baseados em Linux
      & Determinismo \emph{versus} flexibilidade; certificação e atualização de software \\
    \bottomrule
  \end{tabular}
\end{table}

\section{CLP, PAC, controlador de segurança e controlador virtual}

O termo ``CLP'' virou guarda-chuva. Vale separar as categorias
(Tabela~\ref{tab:categorias-controladores}).

\begin{table}[H]
  \centering
  \caption{Categorias de controlador que convivem em uma planta hoje.}
  \label{tab:categorias-controladores}
  \begin{tabular}{p{3.4cm} p{6.4cm} p{4.8cm}}
    \toprule
    Categoria & Característica central & Onde aparece \\
    \midrule
    CLP clássico
      & Lógica discreta e sequencial com ciclo de varredura determinístico; Ladder como linguagem natural
      & Máquinas, linhas de montagem, utilidades \\
    PAC (\emph{programmable automation controller})
      & Mesma base do CLP, com mais memória e processamento, multitarefa real, tipagem forte e forte
        integração em rede
      & Processos que misturam lógica, movimento e malhas PID \\
    Controlador de segurança
      & Adequado a funções críticas de segurança, com nível de integridade (SIL) e conjunto de
        instruções específico para parada de emergência e cortinas de luz
      & Zonas de risco, células robotizadas \\
    CLP virtual ou \emph{soft PLC}
      & Função de controlador executada em hardware de propósito geral, em máquina virtual ou
        contêiner
      & Simulação, ensino, controle remoto, arquiteturas em nuvem \\
    RTU (\emph{remote terminal unit})
      & Controlador rústico, de baixo consumo, cuja função principal é gerir o enlace de
        comunicação, com carimbo de tempo dos eventos
      & Sistemas a distância: saneamento, energia, dutos \\
    \bottomrule
  \end{tabular}
\end{table}

\begin{atencao}
Controlador de segurança não é ``um CLP mais caro''. É um equipamento certificado para uma função
de segurança, com arquitetura e diagnóstico próprios. Lógica de segurança implementada em CLP
comum, sem certificação, \textbf{não} cumpre função de segurança.
\end{atencao}

\section{O que não mudou}

Apesar de meio século de evolução, três decisões de 1969 continuam governando o projeto de
qualquer controlador atual:

\begin{enumerate}
  \item \textbf{o ciclo de varredura} --- ler entradas, processar, escrever saídas, repetir;
  \item \textbf{a linguagem gráfica próxima do esquema elétrico}, que sobreviveu como Ladder
        justamente por ser legível para quem vem do comando elétrico;
  \item \textbf{o foco em diagnóstico} --- sinalizar estado e falha, que era o requisito 3 e segue
        sendo o principal diferencial entre um controlador industrial e um computador comum.
\end{enumerate}

\begin{pratica}
Compare um CLP antigo ainda em operação na sua região (ou a documentação dele) com um controlador
atual. Monte uma tabela com três colunas: \emph{o que era}, \emph{o que é hoje} e \emph{o que
continua igual}. Use a Tabela~\ref{tab:geracoes} como ponto de partida, mas registre o que
encontrar de concreto.
\end{pratica}

\begin{exercicios}
\begin{enumerate}[label=\alph*)]
  \item Por que a indústria automobilística, e não a química ou a de energia, foi a primeira a
        exigir o controlador programável? Relacione a resposta com a frequência de mudança do
        processo.
  \item Escolha três dos dez requisitos da Tabela~\ref{tab:requisitos-originais} e mostre como cada
        um é atendido por um CLP atual.
  \item Explique a diferença entre CLP e PAC sem usar a palavra ``potência''.
  \item Um colega propõe usar um CLP comum, com lógica cuidadosamente escrita, para substituir um
        relé de segurança de uma prensa. Escreva a resposta técnica, em cinco linhas, explicando
        por que a proposta não é aceitável.
  \item Explique o que um CLP virtual permite fazer que o hardware físico não permite, e qual
        problema ele \emph{não} resolve.
\end{enumerate}
\end{exercicios}
